% ============================================================
\chapter{Sober Spaces and Locales}
\label{ch:sober_locales}
% ============================================================

What if you could do topology without ever mentioning points? The idea sounds absurd: since Cantor and Hausdorff, a topological space is a set of \emph{points} equipped with a collection of open sets. Yet starting in the 1960s, a tradition rooted in logic and theoretical computer science---championed by mathematicians like John Isbell and the category theorists---showed that the open sets themselves contain all the information. The lattice of opens, viewed as an algebraic structure, suffices to reconstruct the topology. Spaces where this reconstruction is faithful are called \emph{sober}; the abstract algebraic structure is a \emph{locale}.

This chapter explores the fascinating passage from the concrete (points) to the abstract (opens as autonomous entities), a journey that reveals deep connections between topology, order theory, and logic.

\begin{intuition}
Sober spaces constitute the most natural class of $T_0$-spaces: they are those where the
topology completely determines the point structure, via the correspondence between points
and irreducible closed sets. Locales formalize the idea of ``pointless topology'' and
provide an algebraic framework for topology.
\end{intuition}

% ------------------------------------------------------------
\section{Irreducible Closed Sets}
% ------------------------------------------------------------

\begin{definition}[Irreducible closed set]
A closed set $F$ of a topological space $X$ is \emph{irreducible} if:
\begin{enumerate}
  \item $F \neq \emptyset$;
  \item if $F = F_1 \cup F_2$ with $F_1, F_2$ closed in $F$, then $F = F_1$ or $F = F_2$.
\end{enumerate}
Equivalently, $F$ is irreducible iff for all opens $U_1, U_2$ of $X$ with
$F \cap U_1 \neq \emptyset$ and $F \cap U_2 \neq \emptyset$, we have
$F \cap U_1 \cap U_2 \neq \emptyset$.
\end{definition}

\begin{proposition}[Properties of irreducible closed sets]
\begin{enumerate}
  \item The closure of a singleton $\overline{\{x\}}$ is always irreducible.
  \item The continuous image of an irreducible set is irreducible.
  \item The closure of an irreducible set is irreducible.
  \item If $F$ is irreducible and $F \cap G \neq \emptyset$ with $G$ open, then $F \cap G$
    is irreducible in $G$.
\end{enumerate}
\end{proposition}

\begin{proof}
(1) If $\overline{\{x\}} = F_1 \cup F_2$ with $F_i$ closed, then $x \in F_1$ or
$x \in F_2$, say $x \in F_1$. Then $\overline{\{x\}} \subseteq F_1$, so
$F_1 = \overline{\{x\}}$.

(2) Let $f\colon X \to Y$ be continuous and $F$ irreducible in $X$. If
$f(F) \subseteq G_1 \cup G_2$ with $G_i$ closed, then $F \subseteq f^{-1}(G_1) \cup
f^{-1}(G_2)$ with $f^{-1}(G_i)$ closed, so $F \subseteq f^{-1}(G_i)$ for some $i$,
giving $f(F) \subseteq G_i$. The closure $\overline{f(F)}$ is irreducible by (3).

(3) If $\overline{F} = G_1 \cup G_2$, then $F \subseteq G_1 \cup G_2$, so
$F = (F \cap G_1) \cup (F \cap G_2)$. By irreducibility, $F \subseteq G_i$ for some $i$,
whence $\overline{F} \subseteq G_i$.
\end{proof}

% ------------------------------------------------------------
\section{Sober Spaces}
% ------------------------------------------------------------

\begin{definition}[Sober space]
A topological space $X$ is \emph{sober} if every irreducible closed set has a unique
generic point: for every irreducible closed $F$, there exists a unique $x \in X$ with
$\overline{\{x\}} = F$.
\end{definition}

\begin{theorem}[Characterization of sober spaces]
A space $X$ is sober if and only if:
\begin{enumerate}
  \item $X$ is $T_0$, and
  \item every irreducible closed set has a generic point.
\end{enumerate}
\end{theorem}

\begin{proof}
$(\Rightarrow)$ Uniqueness of the generic point implies $T_0$: if
$\overline{\{x\}} = \overline{\{y\}}$, then $x$ and $y$ are both generic points of the
same irreducible closed set, so $x = y$.

$(\Leftarrow)$ It suffices to show uniqueness. If $\overline{\{x\}} = \overline{\{y\}} = F$,
then since $X$ is $T_0$, $x = y$.
\end{proof}

\begin{example}[Examples of sober spaces]
\begin{enumerate}
  \item Every Hausdorff space is sober (irreducible closed sets are singletons in $T_2$
    spaces).
  \item The spectrum $\mathrm{Spec}(A)$ of a commutative ring is sober (a classical result
    in algebraic geometry).
  \item The Sierpi\'nski space $\mathbb{S}$ is sober.
  \item Every complete lattice with the Scott topology is sober.
\end{enumerate}
\end{example}

\begin{example}[A $T_0$-space that is not sober]
Let $X = \N$ with the topology whose opens are $\emptyset$ and the cofinal sets
$\{n, n+1, n+2, \ldots\}$. The specialization order is the usual order on $\N$. The closed
set $X$ itself is irreducible (any two nonempty opens intersect). But
$X = \overline{\{n\}}$ is impossible for any $n$, since $\overline{\{n\}} = \{0,\ldots,n\}$.
So $X$ is not sober.
\end{example}

% ------------------------------------------------------------
\section{Sobrification}
% ------------------------------------------------------------

\begin{definition}[Sobrification]
Let $X$ be a topological space. The space of \emph{irreducible closed sets} of $X$,
denoted $X^s$ (or $\mathrm{Sob}(X)$), is the set of irreducible closed subsets of $X$,
equipped with the topology whose opens are:
\[
  \tilde{U} = \{F \in X^s : F \cap U \neq \emptyset\}, \quad U \in \mathcal{O}(X).
\]
The map $\eta_X\colon X \to X^s$ defined by $\eta_X(x) = \overline{\{x\}}$ is continuous.
\end{definition}

\begin{theorem}[Properties of sobrification]
\begin{enumerate}
  \item $X^s$ is sober.
  \item The map $U \mapsto \tilde{U}$ is a frame isomorphism
    $\mathcal{O}(X) \cong \mathcal{O}(X^s)$.
  \item $\eta_X$ is universal: for every continuous map $f\colon X \to Y$ with $Y$ sober,
    there exists a unique continuous map $\tilde{f}\colon X^s \to Y$ with
    $f = \tilde{f} \circ \eta_X$.
  \item The sobrification functor $\mathrm{Sob}\colon \mathbf{Top} \to \mathbf{Sob}$ is
    left adjoint to the inclusion $\mathbf{Sob} \hookrightarrow \mathbf{Top}$.
\end{enumerate}
\end{theorem}

\begin{proof}
(1) We show $X^s$ is $T_0$. If $F_1 \neq F_2$ are irreducible closed sets, there exists
an open $U$ with, say, $F_1 \cap U \neq \emptyset$ and $F_2 \cap U = \emptyset$. Then
$F_1 \in \tilde{U}$ and $F_2 \notin \tilde{U}$.

Every irreducible closed set of $X^s$ has a generic point. Let $\mathcal{F}$ be an
irreducible closed subset of $X^s$. Set $F_0 = \overline{\bigcup_{F \in \mathcal{F}} F}$
in $X$. One checks that $F_0$ is irreducible in $X$ and
$\overline{\{F_0\}}^{X^s} = \mathcal{F}$.

(2) The map $U \mapsto \tilde{U}$ preserves arbitrary unions and finite intersections,
and is injective since $\tilde{U_1} = \tilde{U_2}$ implies, for every irreducible closed
$F$, that $F \cap U_1 \neq \emptyset \Leftrightarrow F \cap U_2 \neq \emptyset$. Since
singletons $\overline{\{x\}}$ are irreducible closed, this gives $U_1 = U_2$.

(3) If $f\colon X \to Y$ is continuous and $Y$ sober, define $\tilde{f}(F) =$ the unique
generic point of $\overline{f(F)}$ in $Y$ (which exists since $f(F)$ is irreducible and
$Y$ is sober). Continuity and uniqueness are verified directly.
\end{proof}

% ------------------------------------------------------------
\section{Frames and Locales}
% ------------------------------------------------------------

\begin{definition}[Frame]
A \emph{frame} is a complete lattice $L$ satisfying the infinite distributivity law:
\[
  a \wedge \bigvee_{i\in I} b_i = \bigvee_{i\in I}(a \wedge b_i)
\]
for all $a \in L$ and every family $(b_i)_{i\in I}$ in $L$. A \emph{frame morphism}
$h\colon L \to M$ preserves arbitrary joins and finite meets (including $\top$ and $\bot$).
\end{definition}

\begin{definition}[Locale]
The category $\mathbf{Loc}$ of \emph{locales} is the opposite category of the category of
frames:
\[
  \mathbf{Loc} = \mathbf{Frm}^{\mathrm{op}}.
\]
An object of $\mathbf{Loc}$ is a frame viewed as a ``pointless space,'' and a morphism of
locales $f\colon L \to M$ is a frame morphism $f^*\colon M \to L$.
\end{definition}

\begin{remark}
The reversal of arrows is motivated by the analogy with topological spaces: a continuous
map $f\colon X \to Y$ induces $f^{-1}\colon \mathcal{O}(Y) \to \mathcal{O}(X)$, which
goes in the opposite direction.
\end{remark}

% ------------------------------------------------------------
\section{The Functor $\mathrm{pt}$ and the Adjunction $\mathcal{O} \dashv \mathrm{pt}$}
% ------------------------------------------------------------

\begin{definition}[Points of a locale]
Let $L$ be a locale (i.e., a frame). A \emph{point} of $L$ is a frame morphism
$p\colon L \to \{0,1\}$ (where $\{0,1\}$ is the two-element frame), or equivalently, a
completely prime element $q \in L$ ($q \neq \top$ and $a \wedge b \leq q \Rightarrow
a \leq q$ or $b \leq q$).

The space of points $\mathrm{pt}(L)$ is the set of points of $L$ with the topology whose
opens are $\Sigma_a = \{p : p(a) = 1\}$ for $a \in L$.
\end{definition}

\begin{theorem}[Adjunction $\mathcal{O} \dashv \mathrm{pt}$]
The functors $\mathcal{O}\colon \mathbf{Top} \to \mathbf{Loc}$ and
$\mathrm{pt}\colon \mathbf{Loc} \to \mathbf{Top}$ form an adjunction:
\[
  \mathcal{O} \dashv \mathrm{pt}.
\]
The unit $\eta_X\colon X \to \mathrm{pt}(\mathcal{O}(X))$ sends $x$ to the point
$\eta_X(x)(U) = \begin{cases} 1 & \text{if } x \in U \\ 0 & \text{otherwise}\end{cases}$.
\end{theorem}

\begin{proof}
For a space $X$ and a locale $L$, we establish:
\[
  \mathrm{Hom}_{\mathbf{Top}}(X, \mathrm{pt}(L)) \cong
  \mathrm{Hom}_{\mathbf{Loc}}(\mathcal{O}(X), L).
\]
The right side is $\mathrm{Hom}_{\mathbf{Frm}}(L, \mathcal{O}(X))$. A continuous map
$f\colon X \to \mathrm{pt}(L)$ induces a frame morphism $L \to \mathcal{O}(X)$ sending
$a \mapsto f^{-1}(\Sigma_a)$. Conversely, a frame morphism $h\colon L \to \mathcal{O}(X)$
induces a continuous map $X \to \mathrm{pt}(L)$ sending $x$ to the point
$a \mapsto [x \in h(a)]$.
\end{proof}

\begin{theorem}[Characterization of sober spaces via the adjunction]
A space $X$ is sober if and only if the unit $\eta_X\colon X \to
\mathrm{pt}(\mathcal{O}(X))$ is a homeomorphism.
\end{theorem}

\begin{proof}
The map $\eta_X$ is injective iff $X$ is $T_0$ (distinct points give different ``points''
of the frame). It is surjective iff every point of $\mathcal{O}(X)$ arises from a point of
$X$, which is equivalent to every irreducible closed set having a generic point. Finally,
$\eta_X$ is always a topological embedding.
\end{proof}

% ------------------------------------------------------------
\section{Spatial Locales}
% ------------------------------------------------------------

\begin{definition}[Spatial locale]
A locale $L$ is \emph{spatial} if the canonical frame morphism
$L \to \mathcal{O}(\mathrm{pt}(L))$ is an isomorphism, i.e., if $L$ has ``enough points.''
\end{definition}

\begin{theorem}
The restrictions of the adjunction $\mathcal{O} \dashv \mathrm{pt}$ induce an equivalence
of categories between the category of sober spaces and the category of spatial locales.
\end{theorem}

\begin{proposition}[Non-spatial locales exist]
There exist frames with no points at all: for example, the frame of regular open subsets
of $\R$ modulo a certain ideal is a nontrivial frame without any points.
\end{proposition}

% ------------------------------------------------------------
\section{Sublocales and Nuclei}
% ------------------------------------------------------------

\begin{definition}[Nucleus]
Let $L$ be a frame. A \emph{nucleus} on $L$ is a map $j\colon L \to L$ such that:
\begin{enumerate}
  \item $a \leq j(a)$ (inflationary);
  \item $j(j(a)) = j(a)$ (idempotent);
  \item $j(a \wedge b) = j(a) \wedge j(b)$ (preserves finite meets).
\end{enumerate}
\end{definition}

\begin{theorem}[Sublocales via nuclei]
Sublocales of a locale $L$ correspond bijectively to nuclei on $L$. If $j$ is a nucleus,
the image $j(L) = \{a \in L : j(a) = a\}$ is a frame (with the same meets as $L$ and
modified joins $\bigvee^j a_i = j(\bigvee a_i)$), and the frame morphism
$L \to j(L)$, $a \mapsto j(a)$, defines a locale embedding $j(L) \hookrightarrow L$.
\end{theorem}

\begin{proof}
We verify that $j(L)$ is a frame. It is ordered by the inherited order from $L$. For
meets: if $a, b \in j(L)$, then $j(a \wedge b) = j(a) \wedge j(b) = a \wedge b$, so
$a \wedge b \in j(L)$. For joins: $j(\bigvee a_i)$ is in $j(L)$ and is the smallest
element of $j(L)$ above all $a_i$. Distributivity is verified using that of $L$.
\end{proof}

% ------------------------------------------------------------
\section{Stone Duality}
% ------------------------------------------------------------

\begin{definition}[Bounded distributive lattice]
A \emph{bounded distributive lattice} is a lattice $D$ with a least element $0$ and a
greatest element $1$, satisfying $a \wedge (b \vee c) = (a \wedge b) \vee (a \wedge c)$.
\end{definition}

\begin{definition}[Spectral space]
A topological space $X$ is \emph{spectral} if it is:
\begin{enumerate}
  \item $T_0$ and sober;
  \item compact;
  \item the set of compact opens is closed under finite intersection and forms a base.
\end{enumerate}
\end{definition}

\begin{theorem}[Stone duality for distributive lattices]
There is an equivalence of categories between the category of bounded distributive lattices
and the category of spectral spaces (with spectral continuous maps). The functor in one
direction sends a lattice $D$ to the space of its prime filters, and in the other, sends a
spectral space $X$ to the lattice of its compact open subsets.
\end{theorem}

% ------------------------------------------------------------
\section{Exercises}
% ------------------------------------------------------------

\begin{exercise}
Show that the Sierpi\'nski space $\mathbb{S}$ is sober and describe its irreducible closed
sets.
\end{exercise}

\begin{exercise}
Show that a space is sober if and only if it is a $d$-space and every irreducible closed
set has a generic point.
\end{exercise}

\begin{exercise}
Let $L$ be a frame. Show that $\mathrm{pt}(L)$ is always a sober space.
\end{exercise}

\begin{exercise}
Verify that the frame $\mathcal{O}(\R)$ is spatial (i.e., that $\R$ with the usual topology
gives a spatial locale).
\end{exercise}

\begin{exercise}[Open and closed sublocales]
Let $L$ be a frame and $a \in L$. Show that:
\begin{enumerate}
  \item The nucleus $j_a(x) = a \vee x$ defines a closed sublocale.
  \item The nucleus $j^a(x) = a \Rightarrow x$ (Heyting implication) defines an open
    sublocale.
\end{enumerate}
\end{exercise}

\begin{exercise}
Show that the sobrification of $\N$ (with cofinal open sets) is homeomorphic to
$\N \cup \{\infty\}$ with the Alexandrov topology of the usual order.
\end{exercise}

\begin{exercise}[Heyting algebra]
Show that every frame is a Heyting algebra: for all $a, b$, there exists a greatest
element $c$ with $a \wedge c \leq b$, denoted $a \Rightarrow b$. Show that
$a \Rightarrow b = \bigvee\{c : a \wedge c \leq b\}$.
\end{exercise}
